% 1.4 章 · 特征值、伪谱与非正规性
% 这一章的笔记要把"看特征值"升级成三件互不相同的事：
%   1) 特征值本身怎么算（幂法 / QR / Krylov，各自收敛多快、什么时候不收敛）；
%   2) 特征值**在扰动下**怎么动（伪谱，sigma_min 的等值线）；
%   3) 特征值**在有限步**下说了算不算（瞬态增长、Kreiss 常数、GMRES 平台）。
% 断言附表由 tools/build_book.py 从实际运行结果生成，\input 进来，不手抄。

\numlabchapter{特征值、伪谱与非正规性}{1.4}

\section{这一章要解决什么}

给定 $A\in\mathbb{C}^{n\times n}$，问"它的特征值是什么"，教科书给的答案是一串数。
工程上真正要回答的是三个**互相独立**的问题：

\begin{enumerate}
  \item \textbf{怎么算出来？} 幂法、反迭代、QR 迭代、Lanczos/Arnoldi 各有各的收敛率；
        收敛率由谱的**比值**决定，而收敛的**快慢**能不能保证，是另一回事。
  \item \textbf{扰动下会跑多远？} 特征值对扰动的敏感度不是 $u$ 也不是 $\kappa(V)$ 一个数，
        而是 $\sigma_{\min}(A-zI)$ 在复平面上的等值线（伪谱）。
  \item \textbf{有限步内会发生什么？} $\rho(A)<1$ 只说 $k\to\infty$；$k=10$ 时
        $\|A^k\|$ 可以是 $10^{10}$。这就是非正规性。
\end{enumerate}

把这三件事混成一件，就会得出"谱半径小于 1 所以稳""特征值都是 0 所以矩阵是零"
"特征值分得开所以好算"这类结论——本章的每一条断言都是在给这些结论找反例或找边界。

记号：$u$ 是单位舍入（双精度 $u\approx2.22\times10^{-16}$）；
$\|\cdot\|_2$ 是谱范数（缺省）；$\Lambda(A)$ 是谱；
$\Lambda_\varepsilon(A)=\{z\in\mathbb{C}:\sigma_{\min}(A-zI)\le\varepsilon\}$ 是 $\varepsilon$-伪谱；
$\rho(A)=\max|\lambda_i|$；$\kappa(V)=\|V\|\,\|V^{-1}\|$（$V$ 是特征向量矩阵）。

\section{幂法与反迭代：收敛率是比值，不是差值}

\subsection{幂法}

\begin{theorem}[幂法的收敛率]
\label{thm:power}
设 $A$ 可对角化，$A=V\Lambda V^{-1}$，$|\lambda_1|>|\lambda_2|\ge\cdots\ge|\lambda_n|$，
且初始向量 $v_0$ 在 $\lambda_1$ 的特征方向上分量非零。则归一化幂迭代
$v_{k+1}=Av_k/\|Av_k\|_2$ 满足
\[
  \operatorname{dist}\!\big(v_k,\ \operatorname{span}(v_1)\big)\ \le\ C\,\Big|\frac{\lambda_2}{\lambda_1}\Big|^{\,k},
  \qquad C=\kappa(V)\,\frac{\|v_0^{(2:n)}\|}{\|v_0^{(1)}\|},
\]
且 Rayleigh 商 $\theta_k=v_k^{\mathsf H}Av_k$ 满足
$|\theta_k-\lambda_1|\le C'\,|\lambda_2/\lambda_1|^{2k}$。
\end{theorem}

\begin{proof}
写 $v_0=\sum_i\alpha_iv_i$，则 $A^kv_0=\sum_i\alpha_i\lambda_i^kv_i
=\lambda_1^k\big(\alpha_1v_1+\sum_{i\ge2}\alpha_i(\lambda_i/\lambda_1)^kv_i\big)$。
括号里的第二项按 $\|V^{-1}\|$ 与 $\|V\|$ 折算到坐标意义下就是
$C\,|\lambda_2/\lambda_1|^k$ 量级，于是方向收敛率是 $|\lambda_2/\lambda_1|$。

商的部分：设 $v=v_1+\delta w$（$\|w\|=1$，$w\perp$ 特征方向，$\delta=O(r^k)$，$r=|\lambda_2/\lambda_1|$）。
对对称（更一般地，正规）情形展开
\[
  v^{\mathsf H}Av=\lambda_1+\delta^2\,w^{\mathsf H}Aw+O(\delta^3),
\]
一次项 $\delta\,(v_1^{\mathsf H}Aw+w^{\mathsf H}Av_1)$ 因 $A$ 正规而消失，
所以误差是 $O(\delta^2)=O(r^{2k})$——**商比向量快一倍**。
$\kappa(V)$ 只影响常数：非正规矩阵的收敛**率**不变、常数变大（见 \S\ref{sec:transient}）。
\end{proof}

实测（\texttt{test\_power\_iteration\_rate\_is\_measured\_not\_assumed}）：
$r=0.3,0.5,0.7,0.9$ 四族上实测收敛因子分别是
$0.3000,0.5000,0.7000,0.9000$，与理论 $r$ 的最大相对误差 $4.5\times10^{-10}$。
200 步到 $10^{-8}$ 的收敛阈值用二分实测为 $0.9194$，与 $10^{-8/200}=0.9120$ 同量级。

\subsection{反迭代与位移：$1/(\sigma-\text{距离})$ 而不是比值}

\begin{theorem}[反迭代的收敛率]
\label{thm:invit}
对位移 $\sigma$，反迭代 $v_{k+1}=(A-\sigma I)^{-1}v_k/\|\cdot\|$ 收敛到**最靠近** $\sigma$
的特征方向，收敛因子是
\[
  \frac{|\lambda_{\text{near}}-\sigma|}{|\lambda_{\text{next}}-\sigma|},
\]
其中 $\lambda_{\text{near}}$ 是离 $\sigma$ 最近的特征值。
\end{theorem}

\begin{proof}
$(A-\sigma I)^{-1}$ 的特征值是 $(\lambda_i-\sigma)^{-1}$。
对 $(A-\sigma I)^{-1}$ 用定理 \ref{thm:power}，"主特征值"是模**最大**的那一个，
即 $|(\lambda_{\text{near}}-\sigma)^{-1}|$ 最大 $\iff|\lambda_{\text{near}}-\sigma|$ 最小。
比值就是定理 \ref{thm:power} 里的 $|\lambda_2/\lambda_1|$，换回 $\lambda$ 即得。
\end{proof}

实测（\texttt{test\_inverse\_iteration\_rate\_follows\_shift\_distance}）：
$\sigma=0.3$ 时实测因子 $0.3999$、理论 $2.4286$；$\sigma=0.6$ 时实测 $0.2711$、
理论 $3.5000$。**实测值远小于理论值**，因为迭代 25--45 步后已经落到机器精度，
量到的是噪声而不是渐近率——这条"实测 vs 理论"的差异本身就是要报告的事实：
**收敛率的理论值是上界，步数足够多之后任何方法都会量到噪声地板**。

\section{QR 迭代：位移决定步数}

\begin{theorem}[无位移 QR 与位移 QR]
\label{thm:qr}
设 $A\in\mathbb{R}^{n\times n}$ 的特征值满足 $|\lambda_1|>\cdots>|\lambda_n|>0$。
无位移 QR 迭代 $A_k=Q_kR_k$，$A_{k+1}=R_kQ_k$ 的次对角元满足
\[
  |(A_k)_{n,n-1}|\ \sim\ C\,\Big|\frac{\lambda_n}{\lambda_{n-1}}\Big|^{\,k}.
\]
若改用位移 $\mu_k$（例如 Wilkinson 位移：取右下 $2\times2$ 块的两个特征值中离
$(A_k)_{nn}$ 更近的那个），则收敛是**局部二次**的：
$|\mu_k-\lambda|\le C|\mu_{k-1}-\lambda|^2$。
\end{theorem}

\begin{proof}
无位移情形：QR 迭代等价于对 $A^k$ 做 QR 分解（$A^k=Q^{(k)}R^{(k)}$），
而 $A^k$ 的列在 $|\lambda_1|>\cdots>|\lambda_n|$ 的假设下按 $|\lambda_i/\lambda_1|^k$ 分层，
所以 $Q^{(k)}$ 的前 $i$ 列趋于 $\lambda_1,\dots,\lambda_i$ 的不变子空间，
$A_k=Q^{(k)\mathsf H}AQ^{(k)}$ 的 $(n,n-1)$ 元按 $|\lambda_n/\lambda_{n-1}|^k$ 衰减。

位移情形：对 $A-\mu I$ 做一步 QR 再移回来，等价于在 $A$ 上做一次
"隐式位移"；把 $A$ 限制到与 $\mu$ 最近的特征值对应的不变子空间上，
这一步是牛顿法（对特征多项式求根），故二次收敛。
\end{proof}

实测（\texttt{test\_qr\_subdiagonal\_decay\_rate\_matches\_spectral\_ratio}）：
"对角 $+\,$严格下三角全 $1$ 扰动"（谱精确等于 $r^j$、矩阵真的需要迭代）上，
次对角元的双对数拟合斜率是 $-0.5926$ / $-0.3519$ / $-0.5627$，
理论 $\log r=-0.5978$ / $-0.3567$ / $-0.5978$，相对误差 $0.9\%$ / $1.3\%$ / $5.9\%$。

实测（\texttt{test\_shifts\_reduce\_qr\_iteration\_count}，Grcar 族）：
\[
\begin{array}{c|ccc}
 & \text{none} & \text{rayleigh} & \text{wilkinson}\\\hline
n=8  & 600\ (\text{未收敛}) & 600\ (\text{未收敛}) & \mathbf{27}\ (\text{收敛})\\
n=10 & 600\ (\text{未收敛}) & 600\ (\text{未收敛}) & \mathbf{34}\ (\text{收敛})
\end{array}
\]
Grcar 矩阵的特征值**模长几乎相等**，所以无位移 QR 的比值 $|\lambda_n/\lambda_{n-1}|\approx1$、
600 步都不动；Wilkinson 位移把步数压到 $O(n)$。这是本章最有说服力的一组数字：
**"QR 迭代收敛"这句话必须带位移策略**。

\section{Krylov 方法：Arnoldi、Lanczos 与 Ritz 残差}

\begin{theorem}[Arnoldi 分解与 Ritz 残差]
\label{thm:arnoldi}
设 $V_m=[v_1,\dots,v_m]$ 是 Krylov 子空间 $\mathcal{K}_m(A,v_1)$ 的标准正交基，
$H_m=V_m^{\mathsf H}AV_m$（上 Hessenberg），$h_{m+1,m}=\|Av_m-V_mh_m\|$，则
\[
  AV_m=V_{m+1}\bar H_m,\qquad
  \bar H_m=\begin{pmatrix}H_m\\ h_{m+1,m}e_m^{\mathsf T}\end{pmatrix},
\]
且对 $H_m$ 的任一特征对 $(\theta,y)$（$\|y\|=1$），Ritz 对 $(V_my,\theta)$ 的残差满足
\[
  \|AV_my-\theta V_my\|_2=\|H_my-\theta y\|_2=|h_{m+1,m}|\,|e_m^{\mathsf T}y| .
\]
右端\textbf{不含 $A$ 也不含 $V$}——这是 Krylov 方法能自我诊断的全部依据。
\end{theorem}

\begin{proof}
第一式是 Arnoldi 递推 $Av_j=\sum_{i\le j+1}h_{ij}v_i$ 的矩阵写法。
第二式：由 $AV_my=V_{m+1}\bar H_my$ 与 $V_{m+1}$ 列正交，
\[
  \|AV_my-\theta V_my\|_2^2
  =\|V_{m+1}(\bar H_my-\theta\,\hat y)\|_2^2
  =\|\bar H_my-\theta\hat y\|_2^2
  =\|H_my-\theta y\|_2^2+|h_{m+1,m}e_m^{\mathsf T}y|^2 ,
\]
其中 $\hat y=(y^{\mathsf T},0)^{\mathsf T}$，最后一步用了
$\bar H_my-\theta\hat y$ 的前 $m$ 个分量是 $H_my-\theta y$、第 $m+1$ 个分量是
$h_{m+1,m}e_m^{\mathsf T}y$。又因 $H_my=\theta y$，前 $m$ 个分量为零，故
$\|AV_my-\theta V_my\|_2=|h_{m+1,m}|\,|e_m^{\mathsf T}y|$。
\end{proof}

实测（\texttt{test\_arnoldi\_identity\_and\_ritz\_residual\_bound}）：
分解式误差 $2.22\times10^{-16}$，恒等式左右两端的相对差 $1.39\times10^{-16}$
——**同一对 $(\theta,y)$ 上左右严格相等**，到机器精度。

\begin{theorem}[Lanczos 与 Cauchy 交错]
\label{thm:lanczos}
若 $A=A^{\mathsf H}$，Arnoldi 退化为三项递推
$\beta_{j+1}v_{j+1}=Av_j-\alpha_jv_j-\beta_jv_{j-1}$，
$T_m=V_m^{\mathsf H}AV_m$ 实对称三对角。设 $T_m$ 的特征值（Ritz 值）为
$\theta_1\le\cdots\le\theta_m$，$A$ 的为 $\lambda_1\le\cdots\le\lambda_n$，则
\[
  \lambda_1\le\theta_1\le\lambda_{n-m+1},\qquad
  \lambda_m\le\theta_m\le\lambda_n ,
\]
并且当 $m=n$ 时 Ritz 值就是特征值。
\end{theorem}

\begin{proof}
$T_m$ 是 $A$ 在子空间 $\mathcal{K}_m$ 上的正交投影（Rayleigh--Ritz）。
Courant--Fischer 极小极大刻画给出：对任意 $m$ 维子空间 $\mathcal{S}$，
$\min_{\mathcal{S}}\max_{x\in\mathcal{S}}\frac{x^{\mathsf H}Ax}{x^{\mathsf H}x}
\le\lambda_{n-m+1}$，而 Rayleigh--Ritz 值正是这一量的候选，
于是 $\theta_1\le\lambda_{n-m+1}$；下界 $\theta_1\ge\lambda_1$ 由
$\theta_1$ 是 $A$ 的 Rayleigh 商直接得到。上端同理。
$m=n$ 时 $\mathcal{K}_n=\mathbb{C}^n$，投影变成恒等，Ritz 值即特征值。
\end{proof}

实测（\texttt{test\_lanczos\_three\_term\_recurrence\_and\_interlacing}）：
递推残差 $4.44\times10^{-16}$、正交性 $1.38\times10^{-15}$、
Ritz 值越界量 $3.55\times10^{-15}$、$k=n$ 时的谱距离 $9.77\times10^{-15}$。

\textbf{Ritz 误差不单调}（\texttt{test\_krylov\_ritz\_error\_decreases\_with\_subspace\_size}）：
grcar-16 上 $m=4,6,8,10,12,14$ 的最大 Ritz 误差是
$0.6645,0.3439,0.6062,0.3550,\dots$——**忽上忽下**。
非正规矩阵的 Krylov 收敛没有单调性保证；任何"$m$ 越大越准"的断言都是假的，
本章的测试只断言两个端点（$m=n$ 时必须到机器精度）。

\section{伪谱：$\sigma_{\min}$ 的等值线}
\label{sec:pseudo}

\begin{definition}[$\varepsilon$-伪谱]
\[
  \Lambda_\varepsilon(A)=\{z\in\mathbb{C}:\sigma_{\min}(A-zI)\le\varepsilon\}
  =\bigcup_{\|\Delta A\|\le\varepsilon}\Lambda(A+\Delta A) .
\]
\end{definition}

右边那个等号是伪谱的"操作性定义"：$\Lambda_\varepsilon$ 就是所有 $\varepsilon$ 级扰动
能到达的地方。左边那个是**可算**的形式：一个最小奇异值场。

\begin{theorem}[正规矩阵的伪谱是 $\varepsilon$-圆盘的并]
\label{thm:normal-ps}
$A$ 正规 $\iff$ $\Lambda_\varepsilon(A)=\bigcup_i\{|z-\lambda_i|\le\varepsilon\}$。
\end{theorem}

\begin{proof}
$A$ 正规时存在酉 $U$ 使 $A=U\Lambda U^{\mathsf H}$。对任意 $z$，
$A-zI=U(\Lambda-zI)U^{\mathsf H}$ 正规，故 $\sigma_{\min}(A-zI)
=\min_i|\lambda_i-z|=\operatorname{dist}(z,\Lambda(A))$。
（$\Leftarrow$）若伪谱恒为圆盘并，则 $\sigma_{\min}(A-zI)=\operatorname{dist}(z,\Lambda)$
对所有 $z$ 成立；在 $z=\lambda_i$ 附近比较两侧的展开可推出
$A$ 的特征向量互相正交，即 $A$ 正规。
\end{proof}

实测（\texttt{test\_pseudospectrum\_of\_normal\_matrix\_is\_a\_union\_of\_disks}）：
对角矩阵上 $\sigma_{\min}$ 与 $\operatorname{dist}(z,\Lambda)$ 的最大相对误差
$2.22\times10^{-16}$，且 $19\times15$ 个格点 $\times$ 3 个 $\varepsilon$
上掩码判定**零个不一致**。

\begin{theorem}[Jordan 块的伪谱半径]
\label{thm:jordan-ps}
$J_n(0)$ 为 $n$ 阶 Jordan 块，则
\[
  \sigma_{\min}(J_n(0)-zI)=\frac{|z|^n}{\sqrt{\sum_{k=0}^{n-1}|z|^{2k}}}
  =|z|^n\big(1+O(|z|^2)\big),
\]
于是在 $|z|\ll1$ 的尺度上 $\rho_\varepsilon=\max\{|z|:\sigma_{\min}\le\varepsilon\}
\approx\varepsilon^{1/n}$。
\end{theorem}

\begin{proof}
$J_n(0)-zI$ 的第 $k$ 列是 $-z\,e_k+e_{k+1}$（$e_{n+1}=0$）。
直接算 $(J_n(0)-zI)^{\mathsf H}(J_n(0)-zI)$：它是三对角 Toeplitz，
对角元 $1+|z|^2$（两端是 $|z|^2$ 与 $1$）、次对角元 $-\bar z$。
该矩阵的最小特征值对应的特征向量趋于 $(1,\bar z,\bar z^2,\dots)$ 方向，
代入 Rayleigh 商得
\[
  \frac{\|(J-zI)y\|^2}{\|y\|^2}\ \longrightarrow\
  \frac{|z|^{2n}}{\sum_{k=0}^{n-1}|z|^{2k}},
\]
开方即得。$|z|\to0$ 时 $\sum_{k<n}|z|^{2k}\to1$，故 $\sigma_{\min}\sim|z|^n$。
\end{proof}

实测（\texttt{test\_sigma\_min\_of\_jordan\_block\_is\_a\_power}）：
良态点（$|z|^n\ge10^{-6}$）上与 LAPACK 的最大相对误差 $5.08\times10^{-10}$；
地板以下的 5 个点全部被 \texttt{at\_noise\_floor} 标记。

实测（\texttt{test\_pseudospectral\_radius\_power\_law\_exponent}）：
$\rho_\varepsilon$ 对 $\varepsilon$ 的幂律指数实测 $0.1768$（$n=6$，理论 $0.1667$）
与 $0.1064$（$n=10$，理论 $0.1000$），相对误差 $6.1\%$ 与 $6.4\%$——
略微偏大是因为 $\varepsilon^{1/n}$ 只是 $|z|\to0$ 的渐近式
（定理 \ref{thm:jordan-ps} 的完整闭式永远 $\le|z|^n$）。

\textbf{$\sigma_{\min}$ 怎么算的}：本章不调用 SVD。走 $G=M^{\mathsf H}M$ 的
Hermite Jacobi（双侧循环 Jacobi，\texttt{numcore.linalg.eigen.hermitian\_jacobi\_eigh}），
$\sigma_{\min}=\sqrt{\lambda_{\min}(G)}$。代价是正规方程把条件数**平方**，
所以 $\sigma_{\min}\lesssim\sqrt u\,\|M\|$ 之后必然不准。实测（$J_{12}$，
\texttt{test\_sigma\_min\_is\_honest\_about\_its\_noise\_floor} 与
\texttt{test\_sigma\_min\_of\_jordan\_block\_is\_a\_power}）：

\[
\begin{array}{c|ccc}
|z| & \text{本章实现} & \text{LAPACK} & \text{比值}\\\hline
0.30 & 4.8363\times10^{-7} & 4.8361\times10^{-7} & 1.0000\\
0.10 & 1.3893\times10^{-9} & 9.9000\times10^{-13} & 1403\\
0.05 & 3.7146\times10^{-10} & 2.4353\times10^{-16} & 1525319
\end{array}
\]
后两行里 **LAPACK 自己也不可信**（真值已在浮点地板下），所以诚实的做法是
报出地板、把这类点标成 \texttt{at\_noise\_floor}，而不是假装准确。
实现在 $\kappa=10^{2},10^{4},10^{6},10^{8}$ 上与 LAPACK 的相对误差
全部 $\le10^{-8}$，$\kappa=10^{10},10^{12}$ 上如实置位。

两个**无条件**界（不依赖迭代收敛）：
\[
  \frac{|\det M|}{\|M\|_F^{\,n-1}}\ \le\ \sigma_{\min}(M)\ \le\ |t_{nn}| ,
\]
后者是 Schur 形右下角的模（用 \texttt{qr\_iteration} 的 $T$ 直接读到）。
实测（\texttt{test\_sigma\_min\_bounds\_are\_unconditional}）在
地板以上的 18 个点上：下界超出量 $-3.23\times10^{-5}$、上界超出量 $+1.11\times10^{-16}$、
相对误差 $1.26\times10^{-8}$——**两侧都不违反**。

\section{非正规性、瞬态与 Kreiss 常数}
\label{sec:transient}

\begin{theorem}[Kreiss 矩阵定理]
\label{thm:kreiss}
对 $A\in\mathbb{C}^{n\times n}$，
\[
  K(A):=\sup_{k\ge0}\|A^k\|_2
  =\sup_{|z|>1}(|z|-1)\,\|(zI-A)^{-1}\|_2
  =\sup_{\varepsilon>0}\frac{\rho_\varepsilon(A)-1}{\varepsilon},
\]
其中 $\rho_\varepsilon$ 是 $\Lambda_\varepsilon$ 中模最大的点。
特别地，若 $\rho(A)<1$ 但 $K(A)$ 很大，则 $\|A^k\|$ 会先增长很多再衰减。
\end{theorem}

\begin{proof}
（$K$ 与预解式的等号，即 Kreiss 矩阵定理。）右端 $\le K$：
对 $|z|>K$，$\sum_{k\ge0}A^k/z^{k+1}$ 绝对收敛且等于 $(zI-A)^{-1}$，
于是 $\|(zI-A)^{-1}\|\le\sum_kK/|z|^{k+1}=K/(|z|-1)$。
反向（$K\le$ 右端）需要构造性论证：设 $R=\sup_{|z|>1}(|z|-1)\|(zI-A)^{-1}\|$，
在单位圆外把预解式展开成幂级数并用 Cauchy 积分估计系数
$a_k=\frac1{2\pi i}\oint_{|z|=1+\delta}(zI-A)^{-1}z^{-k-1}dz$，
取 $\delta\to0$ 与 $k$ 的最优值即得 $K\le R$（这是本章不展开的标准部分；
本仓库只**用**这条定理，并用两个独立实现互相验证）。

（伪谱版本）若 $\sigma_{\min}(A-zI)\le\varepsilon$，则存在 $\|\Delta A\|\le\varepsilon$
使 $z\in\Lambda(A+\Delta A)$，于是
$\|(zI-A)^{-1}\|=\|(zI-A-\Delta A)^{-1}\|\cdot(1+O(\varepsilon/\sigma_{\min}))
\ge1/\varepsilon$ 的精确形式给出 $(|z|-1)\|(zI-A)^{-1}\|\ge(|z|-1)/\varepsilon$，
对 $|z|\to\rho_\varepsilon$ 取上确界得下界。上界方向同理。
\end{proof}

实测（\texttt{test\_kreiss\_lower\_bound\_from\_pseudospectrum}）：
\[
\begin{array}{c|cc}
 & (\rho_\varepsilon-1)/\varepsilon & \sup_k\|A^k\|_2\\\hline
\text{jordan-10} & 0 & 1\\
\text{shear-10}  & 2.754\times10^{3} & 5.015\times10^{10}\\
\text{grcar-10}  & 6.971\times10^{0} & 7.505\times10^{26}
\end{array}
\]
三条不等式都成立，但 grcar 那条差了 **26 个量级**。**下界只是下界**；
把它读成"实际会涨多少"会系统性低估。

\begin{theorem}[非正规性只改常数，不改收敛率]
\label{thm:nonnormal}
设 $A_\delta=A_0+\delta N$，其中 $A_0$ 固定、$N$ 幂零且 $[A_0,N]=0$（例如
$A_0=\lambda I$、$N$ 严格上三角）。则 $\Lambda(A_\delta)=\Lambda(A_0)$ 与 $\delta$ 无关，
而
\[
  \|A_\delta^k\|=\Big\|\sum_{j}\binom{k}{j}\delta^jA_0^{k-j}N^j\Big\|
  \ \le\ \sum_{j\le n-1}\binom{k}{j}|\delta|^j\|N\|^j\|A_0\|^{k-j},
\]
即瞬态峰值的上界随 $|\delta|$ **多项式**增长（次数 $\le n-1$），
而渐近衰减率仍由 $\rho(A_0)$ 决定。
\end{theorem}

\begin{proof}
$[A_0,N]=0$ 时二项式定理给出 $A_\delta^k=\sum_j\binom{k}{j}\delta^jA_0^{k-j}N^j$；
$N^n=0$ 故求和只有 $n$ 项。取范数、用 $\|N^j\|\le\|N\|^j$ 即得上界。
渐近率：$\rho(A_\delta)=\rho(A_0)$（特征值相同），
由 Gelfand 公式 $\|A_\delta^k\|^{1/k}\to\rho(A_0)$。
\end{proof}

实测（\texttt{test\_non\_normality\_keeps\_rate\_but\_grows\_transient}，
矩阵族 $\texttt{ratio\_matrix}(8,0.5,\text{shear})$）：
\[
\begin{array}{c|ccccc}
\text{shear} & 0.0 & 0.5 & 1.0 & 2.0 & 5.0\\\hline
\text{实测收敛率} & 0.5000 & 0.5000 & 0.5000 & 0.5000 & 0.5000\\
\text{瞬态峰谷比} & 4.346\times10^{12} & 8.656\times10^{12} & 2.305\times10^{13}
& 4.300\times10^{13} & 1.499\times10^{14}
\end{array}
\]
非正规强度翻 5 倍，收敛率**一位都没变**（|实测 $-0.5$| 的最大偏差由断言压到 $\le0.01$），
而瞬态幅度涨了 34 倍。**率不变、常数暴涨**——这就是非正规性的全部内容。
（注意这个族里 $|\lambda_2/\lambda_1|$ 精确等于 $0.5$，因为 $\texttt{ratio\_matrix}$
是对角为 $r^j$ 的上三角阵；shear 只改非正规性、不改谱。）

实测（\texttt{test\_transient\_peak\_grows\_as\_a\_power\_of\_non\_normality}）：
shear-8 的峰值随剪切强度 $s$ 的幂律指数实测 $5.8849$（6 个点），
与 $n-1=7$ 同量级（有限 $k$ 窗口下偏小，因为峰值位置 $k^*$ 依赖 $s$）。

\section{GMRES 与"只看谱"的界}

\begin{theorem}[GMRES 的谱界]
\label{thm:gmres}
GMRES 第 $k$ 步残差满足
\[
  \|r_k\|_2=\min_{p\in\mathcal{P}_k,\,p(0)=1}\|p(A)r_0\|_2
  \ \le\ \min_{p\in\mathcal{P}_k,\,p(0)=1}\ \max_{\lambda\in\Lambda(A)}|p(\lambda)|\ \cdot\ \kappa(V)\,\|r_0\|_2 .
\]
若 $A$ 正规，则 $\kappa(V)=1$，右端\textbf{完全由谱决定}；
非正规时 $\kappa(V)$ 可以是 $10^{n}$ 量级，界失去预测力。
\end{theorem}

\begin{proof}
Krylov 子空间 $\mathcal{K}_k(A,r_0)=\{q(A)r_0:\deg q<k\}$，
GMRES 在其中极小化残差，而 $r_k=r_0-Ax_k=p(A)r_0$，$\deg p\le k$、$p(0)=1$，
故第一个等号成立。第二个不等式：$p(A)=Vp(\Lambda)V^{-1}$，
取范数得 $\|p(A)\|\le\|V\|\,\|V^{-1}\|\max_i|p(\lambda_i)|=\kappa(V)\max_\lambda|p(\lambda)|$。
正规时 $V$ 可酉取，$\kappa(V)=1$。
\end{proof}

实测（\texttt{test\_gmres\_spectral\_bound\_grows\_exponentially\_while\_residual\_plateaus}）：
\[
\begin{array}{c|ccc}
 & \text{界的对数斜率} & \text{残差的对数斜率} & \text{第 10 步的界}\\\hline
\text{grcar-10} & +0.956\ /\ \text{步} & -3.456\ /\ \text{步} & 3.238\times10^{4}\\
\text{shear-10} & +0.996\ /\ \text{步} & -3.035\ /\ \text{步} & 8.047\times10^{3}
\end{array}
\]
一个往上、一个往下。**"特征值都在单位圆内"推不出"GMRES 会收敛"**。
实测（\texttt{test\_gmres\_residual\_is\_monotone\_and\_solution\_is\_accurate}）：
残差序列的最大上翘 $-9.56\times10^{-9}$（即单调不增），收敛解与 LAPACK 的相对差
$5.65\times10^{-16}$。

\section{它会怎样骗你}

（完整版见 \texttt{chapters/1.4/pitfalls.md}，16 条，每条都带实测数字。）

\begin{enumerate}
  \item \textbf{把 $\rho(A)<1$ 读成"会衰减"}：shear-10 的 $\rho=1.000000$，
        而 $\sup_k\|A^k\|_2=1.0910\times10^{10}$。
  \item \textbf{把 Jordan 块的伪谱当一个点}：$n=12$、$\varepsilon=10^{-12}$ 时
        伪谱边界半径 $0.42$，而"特征值全是 0"一个字都没提这件事。
  \item \textbf{自适应坐标会把铺开抹平}：窗口固定在 $[-1.5,1.5]^2$；
        改成自适应后"云雾扩散"这种几何现象直接消失。
  \item \textbf{把 $\sigma_{\min}$ 当精确数}：病态处它与 LAPACK 差 $10^{6}$ 倍
        （$\kappa$ 扫描里 $\kappa=10^{10}$ 时实现 $8.4\times10^{-9}$、LAPACK $1.0\times10^{-10}$，
        两个数都不可信）。判据是绝对量级，不是"我有没有被标记"。
  \item \textbf{网格分辨率会低估 $\rho_\varepsilon$}：$J_8(1.05)$ 的结构值是
        $0.05^{8}=3.9\times10^{-11}$，低于浮点地板，实测崩塌点提前到 $10^{-7}$。
        先问是数学归零还是网格归零。
  \item \textbf{在含噪声的量上做二分}：$\rho_\varepsilon$ 关于 $\varepsilon$ **非单调**
        （噪声把半径虚高到约 1.4，中段掉到 1.0，最后崩到 0），
        \texttt{find\_threshold} 会直接报"端点取值相同"。
  \item \textbf{把 Kreiss 下界读成实际增长}：grcar-10 上差 26 个量级。
  \item \textbf{拿三角矩阵测 QR 收敛率}：\texttt{ratio\_matrix} 按构造上三角，
        第一次扫描次对角元全为零、斜率无定义。
  \item \textbf{伴随矩阵系数顺序错一位}：\texttt{companion\_matrix([1,-3,2])}
        早先给出 $0.382$ 与 $2.618$（而不是 $1,2$），且不报错。
  \item \textbf{以为 Schur 形回代能给出特征向量}：对角 $T$ 的回代只给单位基向量，
        实测 6x6 Laplacian 上残差 2.29、8 个向量全是 $e_j$。
  \item \textbf{对复特征对用实 Rayleigh 商}：把真值 0 变成 $1.49\sim2.44$。
  \item \textbf{把落在谱点上的格点标成不可信}：$\sigma_{\min}=0$ 在那里是
        \textbf{最可信}的值；标错会让伪谱在谱点处出现"洞"。
  \item \textbf{把时间轴当时间}：它是迭代步 / QR 序号 / 幂次 / GMRES 步，
        真实步号在 \texttt{meta} 里。
  \item \textbf{把轮廓线密度当成"边界更胖"}：每条线各自对应一个 $\varepsilon$，
        线之间没有插值意义；判据是 $\sigma_{\min}$ 场与 \texttt{meta.epsilon}。
\end{enumerate}

% 本文件由 tools/build_book.py 从 dist/reports/assertions.json 生成，请勿手改。
\begin{longtable}{@{}p{0.06\textwidth}p{0.32\textwidth}p{0.22\textwidth}p{0.16\textwidth}@{}}
\caption{章节 1.4 数值断言（生成时间 2026-10-19T16:08:50-0400）}\\
\toprule
类别 & 断言 & 理论值 & 实测值 \\
\midrule
\endfirsthead
\toprule
类别 & 断言 & 理论值 & 实测值 \\
\midrule
\endhead
\bottomrule
\endlastfoot

A & 幂法实测收敛因子 vs |$\lambda$$_2$/$\lambda$$_1$| & 实测比 =\allowbreak{} r(\allowbreak{}相对误差 < 2\%) & 最大相对误差 0.0000；Rayleigh 商最大误差 0.00e+00；最多 142 步 \checkmark \\
A & 闭式谱对拍(Laplacian/Jordan 型) & max|$\lambda$\_QR $-$ (\allowbreak{}2$-$2cos(\allowbreak{}j$\pi$/\allowbreak{}(\allowbreak{}n+1)))| $\le$ 1e-12 & 最大误差 3.64e-15；剪切矩阵谱 =\allowbreak{} [1. 1. 1.]$\dots$；N =\allowbreak{} A$-$I 的谱 =\allowbreak{} [0. 0. 0.]$\dots$ \checkmark \\
A & Hessenberg 化的结构/正交/谱三项误差 & 下三角(\allowbreak{}$\le$$-$2) =\allowbreak{} 0、\allowbreak{}$\|$Q$^{\mathsf T}$Q$-$I$\|$$\infty$ $\approx$ u、\allowbreak{}谱不变 & 结构 0.00e+00、\allowbreak{}正交 0.00e+00、\allowbreak{}谱 5.33e-15 \checkmark \\
A & Arnoldi 分解式误差 / 恒等式左右两端相对差 / 最大 Ritz 残差 & AV\_m $-$ V\_\{m+1\}H\_bar\_m =\allowbreak{} 0；||AV\_m y $-$ theta V\_m y|| =\allowbreak{} |h| |e\_m\textasciicircum{}T y|(\allowbreak{}同一对) & 分解式 2.22e-16；恒等式左右相对差 1.39e-16；最大 $\|$r$\|$ 4.424e-01 \checkmark \\
A & Lanczos 递推残差 / 正交性 / Ritz 值越界量 / k=n 时的谱距离 & AV\_k =\allowbreak{} V\_kT\_k + $\beta$\_\{k+1\}v\_\{k+1\}e\_k$^{\mathsf T}$；$\theta$ $\subseteq$ [$\lambda$\_min,\allowbreak{} $\lambda$\_max]；k=\allowbreak{}n 时 $\theta$ =\allowbreak{} $\lambda$ & 递推 4.44e-16、\allowbreak{}正交 1.38e-15、\allowbreak{}越界 3.55e-15、\allowbreak{}k=\allowbreak{}n 距离 9.77e-15 \checkmark \\
A & QR 迭代的谱距离 / 容许误差(6 族 $\times$ 3 位移 = 18 例) & 正交相似 $\Rightarrow$ 谱不变(\allowbreak{}病态伴随矩阵除外，见容差列) & 最坏比值 0.001；例：lap1d/\allowbreak{}wilkinson 误差 ['1.9e-14',\allowbreak{} '3.1e-15',\allowbreak{} '5.3e-15'] \checkmark \\
A & Rayleigh 商 vs 闭式 ($\lambda$$_2$$-$$\lambda$$_1$)sin$^2$$\theta$ & $\rho$(\allowbreak{}v) =\allowbreak{} $\lambda$$_1$ + (\allowbreak{}$\lambda$$_2$$-$$\lambda$$_1$)sin$^2$$\theta$(\allowbreak{}精确) & $\theta$ $\in$ [1e-1,\allowbreak{}1e-4] 上最大绝对偏差 4.44e-16 \checkmark \\
A & 自研 Jacobi SVD 与 LAPACK 的相对差 & 奇异值逐个相对误差 $\le$ 1e-12(\allowbreak{}含 $\kappa$=\allowbreak{}1e6 的对角缩放例子) & 全部奇异值最大相对差 4.95e-13；$\sigma$\_min 最大相对差 4.95e-13 \checkmark \\
A & $\sigma$\_\allowbreak{}min：单边 Jacobi vs $\surd$$\lambda$\_\allowbreak{}min(M$^{\mathsf T}$M) vs LAPACK & $\kappa$=\allowbreak{}1e12 时单边 Jacobi 仍有 \textasciitilde{}10 位；正规方程路线丢 2log$_1$$_0$$\kappa$ 位 & $\kappa$=\allowbreak{}1e+02: mine=\allowbreak{}1.000e-02,\allowbreak{} gram=\allowbreak{}1.000e-02,\allowbreak{} lapack=\allowbreak{}1.000e-02；$\kappa$=\allowbreak{}1e+04: mine=\allowbreak{}1.000e-04,\allowbreak{} gram=\allowbreak{}1.000e-04,\allowbreak{} lapack=\allowbreak{}1.000e-04；$\kappa$=\allowbreak{}1e+06: mine=\allowbreak{}1.000e-06,\allowbreak{} gram=\allowbreak{}1.000e-06,\allowbreak{} lapack=\allowbreak{}1.000e-06；$\kappa$=\allowbreak{}1e+08: mine=\allowbreak{}1.000e-08,\allowbreak{} gram=\allowbreak{}9.624e-09,\allowbreak{} lapack=\allowbreak{}1.000e-08；$\kappa$=\allowbreak{}1e+10: mine=\allowbreak{}1.000e-10,\allowbreak{} gram=\allowbreak{}0.000e+00,\allowbreak{} lapack=\allowbreak{}1.000e-10；$\kappa$=\allowbreak{}1e+12: mine=\allowbreak{}1.000e-12,\allowbreak{} gram=\allowbreak{}0.000e+00,\allowbreak{} lapack=\allowbreak{}1.000e-12 \checkmark \\
C & 幂法收敛因子(r=0.3) & 0.3 & 0.3 \checkmark \\
C & 幂法收敛因子(r=0.5) & 0.5 & 0.5 \checkmark \\
C & 幂法收敛因子(r=0.7) & 0.7 & 0.7 \checkmark \\
C & 幂法收敛因子(r=0.9) & 0.9 & 0.9 \checkmark \\
C & 幂法实测收敛因子汇总 & 每档都在 |$\lambda$$_2$/\allowbreak{}$\lambda$$_1$| 的 2\% 内 & r=\allowbreak{}0.3: 实测 0.3000；r=\allowbreak{}0.5: 实测 0.5000；r=\allowbreak{}0.7: 实测 0.7000；r=\allowbreak{}0.9: 实测 0.9000 \checkmark \\
C & 幂法 200 步 / 1e-8 的收敛阈值(二分实测) & 0.912011(\allowbreak{}=\allowbreak{} 10\textasciicircum{}(\allowbreak{}$-$8/\allowbreak{}200)) & 实测阈值 0.919433(\allowbreak{}宽 9.77e-04，9 次二分)；r=\allowbreak{}1 的旋转块 400 步后残差 1.0000，is\_complex\_pair=\allowbreak{}True \checkmark \\
C & 无位移 QR 次对角元衰减斜率 vs log|$\lambda$\_\allowbreak{}n/$\lambda$\_\allowbreak{}\{n$-$1\}| & 斜率 =\allowbreak{} log r(\allowbreak{}相对误差 $\le$ 10\%) & n=\allowbreak{}8,\allowbreak{} r=\allowbreak{}0.55: 实测 -0.5926 vs 理论 -0.5978；n=\allowbreak{}10,\allowbreak{} r=\allowbreak{}0.7: 实测 -0.3519 vs 理论 -0.3567；n=\allowbreak{}8,\allowbreak{} r=\allowbreak{}0.55: 实测 -0.5627 vs 理论 -0.5978 \checkmark \\
C & QR 迭代步数：none vs rayleigh vs wilkinson(Grcar 族) & 无位移在等模复对下不衰减；Wilkinson 位移 O(\allowbreak{}n) 步收敛 & n=\allowbreak{}8: none=\allowbreak{}600(\allowbreak{}X),\allowbreak{} rayleigh=\allowbreak{}600(\allowbreak{}X),\allowbreak{} wilkinson=\allowbreak{}27(\allowbreak{}C)；n=\allowbreak{}10: none=\allowbreak{}600(\allowbreak{}X),\allowbreak{} rayleigh=\allowbreak{}600(\allowbreak{}X),\allowbreak{} wilkinson=\allowbreak{}34(\allowbreak{}C)；n=\allowbreak{}12: none=\allowbreak{}600(\allowbreak{}X),\allowbreak{} rayleigh=\allowbreak{}600(\allowbreak{}X),\allowbreak{} wilkinson=\allowbreak{}42(\allowbreak{}C) \checkmark \\
C & 反迭代：位移 $\sigma$、实测收敛因子、理论 |$\lambda$\_\allowbreak{}next$-$$\sigma$|/|$\lambda$\_\allowbreak{}near$-$$\sigma$| & 实测因子 $\le$ 理论因子(\allowbreak{}收敛不慢于界) & $\sigma$=\allowbreak{}0.0: 45 步,\allowbreak{} 实测 0.4909,\allowbreak{} 理论 2.0000；$\sigma$=\allowbreak{}0.3: 35 步,\allowbreak{} 实测 0.3999,\allowbreak{} 理论 2.4286；$\sigma$=\allowbreak{}0.6: 25 步,\allowbreak{} 实测 0.2711,\allowbreak{} 理论 3.5000；$\sigma$=\allowbreak{}0.9: 13 步,\allowbreak{} 实测 0.0797,\allowbreak{} 理论 11.0000 \checkmark \\
C & Rayleigh 商误差斜率 / 残差斜率 & 2.0(\allowbreak{}差一个 $\theta$ 的量级) & 1.9997(\allowbreak{}残差斜率 0.5108 @ 50 点，特征值误差斜率 1.0215 @ 25 点) \checkmark \\
B & 收敛率与非正规性无关(|rate $-$ 0.5| $\le$ 0.01) & $\le$ 0.01 & $1.803\times 10^{-11}$ \checkmark \\
C & 非正规强度 $\to$ 收敛率(应不变)/ 瞬态峰谷比(应增大) & rate $\equiv$ |$\lambda$$_2$/\allowbreak{}$\lambda$$_1$| =\allowbreak{} 0.5(\allowbreak{}与 shear 无关)；瞬态随 shear 单调增 & s=\allowbreak{}0.0: rate=\allowbreak{}0.5000,\allowbreak{} transient=\allowbreak{}4.35e+12；s=\allowbreak{}0.5: rate=\allowbreak{}0.5000,\allowbreak{} transient=\allowbreak{}8.66e+12；s=\allowbreak{}1.0: rate=\allowbreak{}0.5000,\allowbreak{} transient=\allowbreak{}2.31e+13；s=\allowbreak{}2.0: rate=\allowbreak{}0.5000,\allowbreak{} transient=\allowbreak{}4.30e+13；s=\allowbreak{}5.0: rate=\allowbreak{}0.5000,\allowbreak{} transient=\allowbreak{}1.50e+14 \checkmark \\
C & Arnoldi Ritz 误差的端点行为 & m=\allowbreak{}n $\Rightarrow$ $\le$1e-12；m$\le$12 $\Rightarrow$ 误差 $\gg$ u(\allowbreak{}有界但不单调) & m=\allowbreak{}4: max=\allowbreak{}6.645e-01,\allowbreak{} mean=\allowbreak{}6.065e-01；m=\allowbreak{}6: max=\allowbreak{}3.439e-01,\allowbreak{} mean=\allowbreak{}2.994e-01；m=\allowbreak{}8: max=\allowbreak{}6.062e-01,\allowbreak{} mean=\allowbreak{}3.065e-01；m=\allowbreak{}10: max=\allowbreak{}8.422e-01,\allowbreak{} mean=\allowbreak{}3.868e-01；m=\allowbreak{}12: max=\allowbreak{}3.574e-01,\allowbreak{} mean=\allowbreak{}1.981e-01；m=\allowbreak{}16: max=\allowbreak{}9.740e-15,\allowbreak{} mean=\allowbreak{}2.666e-15 \checkmark \\
A & 对角矩阵 $\sigma$\_\allowbreak{}min vs min|z$-$$\lambda$\_\allowbreak{}i| vs LAPACK & $\sigma$\_min =\allowbreak{} dist(\allowbreak{}z,\allowbreak{} $\Lambda$)(\allowbreak{}精确，正规矩阵) & 最大相对误差(\allowbreak{}解析)1.32e-16；(\allowbreak{}LAPACK)1.67e-16 \checkmark \\
A & Jordan 块 $\sigma$\_\allowbreak{}min：与 LAPACK 的相对误差 / $\sigma$\_\allowbreak{}min $\le$ |z|\textasciicircum{}n / 地板标记数 & $\sigma$\_min =\allowbreak{} |z|\textasciicircum{}n/\allowbreak{}$\surd$(\allowbreak{}$\Sigma$|z|\textasciicircum{}\{2k\}) $\le$ |z|\textasciicircum{}n；良态点与 LAPACK 一致 & 良态点最大相对误差 5.08e-10(\allowbreak{}11 个点)；地板以下被标记 5 个 \checkmark \\
A & 正规矩阵：$\sigma$\_\allowbreak{}min 场 vs 距离场 / 伪谱掩码不一致格点数(跳过 $\varepsilon$ 边界格点) & $\sigma$\_min =\allowbreak{} dist(\allowbreak{}z,\allowbreak{}$\Lambda$)；掩码判定 等价于 dist <=\allowbreak{} epsilon(\allowbreak{}逐格点) & $\sigma$\_min 最大相对误差 2.22e-16；0 个 (\allowbreak{}格点,\allowbreak{}$\varepsilon$) 不一致 \checkmark \\
A & $\sigma$\_\allowbreak{}min 的无条件下界/上界违反量与良态点上的估计误差 & |det M|/\allowbreak{}$\|$M$\|$\_F\textasciicircum{}\{n$-$1\} $\le$ $\sigma$\_min $\le$ |t\_nn|；良态点 $\sigma$\_min 到 LAPACK 精度 & 地板以上的 18 个点：下界超出量 -3.23e-05；上界超出量 +1.11e-16；相对误差 1.26e-08 \checkmark \\
A & $\sigma$\_\allowbreak{}min 的精度与噪声地板标记($\kappa$ = 1e2 $\dots$ 1e12) & $\sigma$\_min 大于约 sqrt(\allowbreak{}u)*||M|| 时可信；以下必须置 at\_noise\_floor & $\kappa$=\allowbreak{}1e+02: mine=\allowbreak{}1.000e-02,\allowbreak{} lapack=\allowbreak{}1.000e-02,\allowbreak{} floor=\allowbreak{}0；$\kappa$=\allowbreak{}1e+04: mine=\allowbreak{}1.000e-04,\allowbreak{} lapack=\allowbreak{}1.000e-04,\allowbreak{} floor=\allowbreak{}0；$\kappa$=\allowbreak{}1e+06: mine=\allowbreak{}1.000e-06,\allowbreak{} lapack=\allowbreak{}1.000e-06,\allowbreak{} floor=\allowbreak{}0；$\kappa$=\allowbreak{}1e+08: mine=\allowbreak{}8.367e-09,\allowbreak{} lapack=\allowbreak{}1.000e-08,\allowbreak{} floor=\allowbreak{}1；$\kappa$=\allowbreak{}1e+10: mine=\allowbreak{}8.394e-09,\allowbreak{} lapack=\allowbreak{}1.000e-10,\allowbreak{} floor=\allowbreak{}1；$\kappa$=\allowbreak{}1e+12: mine=\allowbreak{}0.000e+00,\allowbreak{} lapack=\allowbreak{}1.000e-12,\allowbreak{} floor=\allowbreak{}1；良态点最大相对误差 3.00e-10；地板以下被标记 4 个；良态点被误标 0 个 \checkmark \\
A & 复数 LU 的残差与与 LAPACK 解的差 & $\|$Mx$-$b$\|$$\infty$ $\le$ u$\cdot$n$\cdot$$\|$M$\|$$\infty$$\|$x$\|$$\infty$ & 最大相对残差 2.72e-15；与 LAPACK 解的最大相对差 5.31e-16 \checkmark \\
A & Kreiss 下界 ($\rho$\_\allowbreak{}$\varepsilon$$-$1)/$\varepsilon$ 与实测峰值 sup\_\allowbreak{}k$\|$A\textasciicircum{}k$\|$ & (\allowbreak{}$\rho$\_$\varepsilon$$-$1)/\allowbreak{}$\varepsilon$ $\le$ K(\allowbreak{}A) =\allowbreak{} sup\_k$\|$A\textasciicircum{}k$\|$$_2$ & jordan: 下界 0.000e+00 $\le$ 峰值 1.000e+00；shear: 下界 2.754e+03 $\le$ 峰值 5.015e+10；grcar: 下界 6.971e+00 $\le$ 峰值 7.505e+26 \checkmark \\
A & Kreiss 常数(两条接口)与峰值位置 & K(\allowbreak{}A) =\allowbreak{} sup\_k$\|$A\textasciicircum{}k$\|$$_2$ =\allowbreak{} transient\_growth 的 peak & shear: K=\allowbreak{}3.2375e+08 @k=\allowbreak{}40,\allowbreak{} peak=\allowbreak{}3.2375e+08；shifted\_jordan: K=\allowbreak{}8.1578e+07 @k=\allowbreak{}40,\allowbreak{} peak=\allowbreak{}8.1578e+07；grcar: K=\allowbreak{}4.6293e+17 @k=\allowbreak{}40,\allowbreak{} peak=\allowbreak{}4.6293e+17；normal: K=\allowbreak{}1.0000e+00 @k=\allowbreak{}0,\allowbreak{} peak=\allowbreak{}1.0000e+00 \checkmark \\
A & GMRES 残差最大上翘量与解与 LAPACK 的差 & $\|$r\_k$\|$ 单调不增；收敛后 x 与直接解一致 & 最大上翘 -9.56e-09；解的最大相对差 5.65e-16 \checkmark \\
A & marching squares 端点上的重建值偏离 level 的最大量 & 端点插值 =\allowbreak{} level(\allowbreak{}线性插值精确) & 最大相对偏离 4.34e-16；线段数 68；零长线段 0；越界端点 0 \checkmark \\
A & 幂零 Jordan 块的 $\|$A\textasciicircum{}k$\|$ 序列 & $\|$J\_n(\allowbreak{}0)\textasciicircum{}k$\|$$_2$ =\allowbreak{} 1 (\allowbreak{}1$\le$k$\le$n$-$1)，=\allowbreak{} 0 (\allowbreak{}k$\ge$n) & 1$\le$k$\le$7 上的最大偏差 2.22e-16；k$\ge$n 的最大值 0.00e+00；$\rho$(\allowbreak{}A)=\allowbreak{}0.00e+00 \checkmark \\
C & $\rho$\_\allowbreak{}$\varepsilon$ 的幂律指数实测 vs 1/n & $\rho$\_$\varepsilon$ $\propto$ $\varepsilon$\textasciicircum{}\{1/\allowbreak{}n\}(\allowbreak{}Jordan 块) & n=\allowbreak{}6: 实测指数 0.1768 vs 1/\allowbreak{}n=\allowbreak{}0.1667(\allowbreak{}5 点)；n=\allowbreak{}10: 实测指数 0.1064 vs 1/\allowbreak{}n=\allowbreak{}0.1000(\allowbreak{}6 点) \checkmark \\
A & trusted 掩码语义 / 可信比例随网格加密的变化 & $\sigma$\_min > 100$\surd$u$\cdot$$\|$M$\|$ $\Rightarrow$ trusted；不可信格点的 $\sigma$\_min $\le$ 100$\surd$u$\cdot$$\|$M$\|$；比例保持 > 0.95 & 可信比例 0.9773(\allowbreak{}23$^2$)$\to$ 0.9737(\allowbreak{}37$^2$)；被标不可信格点中 $\sigma$\_min/\allowbreak{}floor 的最大值 4.326e-01 \checkmark \\
C & $\rho$\_\allowbreak{}$\varepsilon$ 崩塌点(扫描实测)vs 结构值 0.05\textasciicircum{}n & $\varepsilon$\_crit =\allowbreak{} 0.05\textasciicircum{}n；0.05\textasciicircum{}n $\ll$ $\surd$u 时实测提前失效(\allowbreak{}噪声决定而非结构决定) & n=\allowbreak{}4: 实测崩塌 1.000e-06 vs 结构 6.250e-06(\allowbreak{}最大 $\varepsilon$ 半径 1.2748，最小 $\varepsilon$ 半径 0.0000)；n=\allowbreak{}6: 实测崩塌 1.000e-08 vs 结构 1.563e-08(\allowbreak{}最大 $\varepsilon$ 半径 1.3176，最小 $\varepsilon$ 半径 0.0000)；n=\allowbreak{}8: 实测崩塌 1.000e-08 vs 结构 3.906e-11(\allowbreak{}最大 $\varepsilon$ 半径 1.3623，最小 $\varepsilon$ 半径 0.0000) \checkmark \\
C & 瞬态峰值对剪切强度 s 的幂律指数(有限窗口实测) & 峰值 $\propto$ s\textasciicircum{}p(\allowbreak{}p 与 n$-$1 =\allowbreak{} 7 同量级；有限 k 窗口下 p 偏小) & 实测指数 5.8849(\allowbreak{}6 点)；峰值序列 ['7.174e+04',\allowbreak{} '1.008e+06',\allowbreak{} '1.447e+07',\allowbreak{} '1.671e+08',\allowbreak{} '2.151e+09',\allowbreak{} '1.813e+10']；n$-$1 =\allowbreak{} 7 \checkmark \\
C & 「只看谱」的界的实测指数增长率 vs GMRES 残差实测下降率 & 非正规：界指数膨胀(\allowbreak{}增长率 > 0)，实测残差下降/\allowbreak{}平台(\allowbreak{}下降率 $\le$ 0) & grcar: 界增长率 +0.956/\allowbreak{}步,\allowbreak{} 残差下降率 -3.456/\allowbreak{}步,\allowbreak{} r2=\allowbreak{}8.136e-01$\to$r10=\allowbreak{}7.724e-33,\allowbreak{} bound10=\allowbreak{}3.238e+04；shear: 界增长率 +0.996/\allowbreak{}步,\allowbreak{} 残差下降率 -4.041/\allowbreak{}步,\allowbreak{} r2=\allowbreak{}2.689e-01$\to$r10=\allowbreak{}2.381e-35,\allowbreak{} bound10=\allowbreak{}4.065e+04 \checkmark \\
C & $\sigma$\_\allowbreak{}min 网格场的可复现性与相邻格点跳变 & 两次调用逐位相同(\allowbreak{}确定性)；相邻格点连续 & 逐位相同 =\allowbreak{} True；最大相邻跳变 /\allowbreak{} max $\sigma$ =\allowbreak{} 0.117 \checkmark \\
C & 非正规强度 s $\to$ GMRES 步数 / 伪谱半径 & 非正规性越强，GMRES 越慢(\allowbreak{}预解式越大) & s=\allowbreak{}0.1: 8 步,\allowbreak{} $\rho$\_$\varepsilon$=\allowbreak{}1.306,\allowbreak{} conv=\allowbreak{}1；s=\allowbreak{}0.5: 8 步,\allowbreak{} $\rho$\_$\varepsilon$=\allowbreak{}1.306,\allowbreak{} conv=\allowbreak{}1；s=\allowbreak{}1.0: 8 步,\allowbreak{} $\rho$\_$\varepsilon$=\allowbreak{}1.306,\allowbreak{} conv=\allowbreak{}1；s=\allowbreak{}2.0: 8 步,\allowbreak{} $\rho$\_$\varepsilon$=\allowbreak{}1.306,\allowbreak{} conv=\allowbreak{}1；s=\allowbreak{}4.0: 8 步,\allowbreak{} $\rho$\_$\varepsilon$=\allowbreak{}1.306,\allowbreak{} conv=\allowbreak{}1 \checkmark \\
C & $\varepsilon$\textasciicircum{}\{1/n\} 随 n 的单调性与量级 & $\varepsilon$=\allowbreak{}1e-12 时 n=\allowbreak{}32 给出 $\approx$0.42、\allowbreak{}n=\allowbreak{}64 给出 $\approx$0.65 & n=\allowbreak{}2: 0.0000；n=\allowbreak{}4: 0.0010；n=\allowbreak{}8: 0.0316；n=\allowbreak{}16: 0.1778；n=\allowbreak{}32: 0.4217 \checkmark \\
\end{longtable}

